\scanpage{42}
ment'' évidente, qui ici dit \deleted{\unsure{qu'}}\inserted{que la donnée}
d'une gerbe \inserted{$\mathcal G_X$}
sur $X$, revient à celle de gerbes $\mathcal G_{X^*}$ sur
$X^*$ et $\mathcal G_{Y,X}$ sur $\mathcal V_{Y,X}$, et d'une
\deleted{iso.}\inserted{équiv.} de \deleted{\unsure{recollem.}} gerbes
\[
(24)\qquad \mathcal G_{Y,X}\mid\mathcal V^*_{Y,X}\simeq
\mathcal G_{X^*}\mid\mathcal V^*_{Y,X}.
\]
Or sur un $K(\pi,1)$, une $A$-gerbe
s'interprète canoniquement comme une
\uline{extension} de $\pi$ par $A$. Donc la donnée
\margin{\diagram{\begin{tikzcd}[column sep=small,row sep=small]
D\arrow[r]\arrow[d]&\Pi_1(Y)\\
\Pi_1(X^*)&
\end{tikzcd}}}
d'une gerbe sur $X\deleted{{}\unsure{\simeq} S^2}$ revient à celle
d'\deleted{\unsure{une ext}} deux extensions, l'une de
$\Pi_1(Y)$ par $A$, l'autre de $\Pi_1(X^*)$ par $A$
et d'un isom. entre les deux extensions
de $D$ par $A$ qui s'en déduisent\unsure{,}
par image inverse par les hom (10).
(cette description marche dès que les
``pièces'' $Y$, $X'$, $\mathcal V^*_{Y,X}$ sont connexes, et
que leur $\Pi_2$ est nul). Dans le cas
% Pale lower-left foliation, another hand; not authorial page numbering.
\margin{41}
